% Part: computability
% Chapter: recursive-functions
% Section: normal-form

\documentclass[../../../include/open-logic-section]{subfiles}

\begin{document}

\olfileid{cmp}{rec}{nft}
\olsection{સામાન્ય સ્વરૂપ પ્રમેય}

\begin{thm}[ક્લીનીનું સામાન્ય સ્વરૂપ પ્રમેય]
\ollabel{thm:kleene-nf}
એવો આદિમ પુનરાવર્તી સંબંધ $T(e, x, s)$ અને એવું આદિમ
પુનરાવર્તી વિધેય $U(s)$ છે કે નીચેનો ગુણધર્મ મળે છે: $f$ કોઈ પણ
એક-આર્ગ્યુમેન્ટવાળું આંશિક પુનરાવર્તી વિધેય હોય, તો કોઈ~$e$ માટે
\[
f(x) \simeq U(\umin{s}{T(e, x, s)})
\]
દરેક $x$ માટે સાચું છે.
\sourcecorrection{OLCMP-010}{સ્થિર અંગ્રેજી મૂળ અહીં ``કોઈ પણ
આંશિક પુનરાવર્તી વિધેય'' કહે છે, પરંતુ પ્રદર્શિત સૂત્ર અને
$T(e,x,s)$ એક ઇનપુટ~$x$ માટે લખાયાં છે. ગુજરાતી વાક્ય આ લખેલા
સ્વરૂપને એક-આર્ગ્યુમેન્ટવાળા વિધેય માટે સ્પષ્ટ કરે છે; અનેક
આર્ગ્યુમેન્ટનું સ્વરૂપ જુદાં ટ્યુપલ-સંકેતીકરણથી મળે છે.}
\end{thm}

\begin{explain}
સામાન્ય સ્વરૂપ પ્રમેયની સાબિતી વિગતભરી છે, પરંતુ મૂળ વિચાર સરળ
છે. દરેક આંશિક પુનરાવર્તી વિધેયને એક \emph{સૂચક}~$e$ હોય છે:
સહજ રીતે, તે તેના કાર્યક્રમ અથવા વ્યાખ્યાનો કોડ કરતી સંખ્યા છે.
$f(x)
\fdefined$ હોય, તો ગણતરી વ્યવસ્થિત રીતે નોંધાઈને કોઈ સંખ્યા~$s$
વડે સંકેતીકૃત થઈ શકે. $s$ એ ઇનપુટ~$x$ પર~$f$ની ગણતરીનો કોડ
છે કે નહીં તે માત્ર $x$ અને વ્યાખ્યાના સૂચક~$e$ વડે આદિમ
પુનરાવર્તી રીતે ચકાસી શકાય છે. પરિણામે સંબંધ~$T$---``સૂચક~$e$
ધરાવતા વિધેયની ઇનપુટ~$x$ માટે ગણતરી છે અને $s$ એ ગણતરીનો
કોડ છે''---આદિમ પુનરાવર્તી છે. ગણતરીની સંપૂર્ણ નોંધ~$s$ મળી જાય
તો~$s$નું ``પરિણામ'' એટલે~$f(x)$નું મૂલ્ય, અને તેને પણ~$s$માંથી
આદિમ પુનરાવર્તી રીતે મેળવી શકાય છે.

સામાન્ય સ્વરૂપ પ્રમેય બતાવે છે કે કોઈ પણ આંશિક પુનરાવર્તી વિધેય
વ્યાખ્યાયિત કરવા અનિયત મર્યાદા સુધીની માત્ર એક શોધ જરૂરી છે.
મૂળભૂત રીતે, સૂચક~$e$ ધરાવતા વિધેયની ઇનપુટ~$x$ માટેની ગણતરીનો
કોડ મળે ત્યાં સુધી બધી સંખ્યાઓમાં શોધીએ છીએ. આંશિક પુનરાવર્તી
વિધેયોનાં ``નામો'' તરીકે સંખ્યાઓ~$e$ વાપરી શકીએ અને પ્રમેયના
સમીકરણથી વ્યાખ્યાયિત વિધેય~$f$ માટે $\cfind{e}$ લખી શકીએ.
નોંધો કે એક જ આંશિક પુનરાવર્તી વિધેયને એકથી વધુ સૂચકો હોઈ શકે
છે---હકીકતમાં, દરેકને અનંત સૂચકો હોય છે.
\end{explain}
\end{document}
